\section{Conclusion}
\label{sec:conclusion}

This software study establishes an algorithmic replay--plasticity trade-off for
memory-constrained, patient-specific adaptation of a tiny ECG transformer, under
a hardware-motivated 16\,KiB persistent-state constraint. The architecture
exhibits a deterministic replay-point cliff---tokens and feed-forward tensors
are several times larger than the raw beat, and only the post-pool
representation is storage-com\-pres\-sive---which makes maximum-volume post-pool
replay the naive budget-optimal design. Over all 21 eligible MIT--BIH DS2
records and five seeds, that naive choice is not established as uniquely optimal:
rank-1 and
rank-2 encoder low-rank adapters with only tens of raw replay samples raise
per-patient \macrof{} from $0.666$ to $0.810$ and $0.811$, improve 16 of 21
patients (versus 12 of 21 for maximum head-only replay), and are both
significant against the frozen model under the pre-registered primary test
(Holm $p\le0.017$) with large effect sizes.
Paired ablations attribute the target gain to encoder plasticity ($p = 0.005$)
but leave replay's direct contribution statistically unsupported ($p = 0.473$),
while replay provides substantially better source retention than the tested
zero-byte $B$-factor anchor under the evaluated configuration; within the
measured 12-record Adam rank-1/rank-2 panel, increasing the adaptation-state
budget from 8 to 64\,KiB produced no statistically detectable benefit, with 14 of
18 fixed-configuration contrasts meeting the pre-specified $\pm 0.05$ equivalence
criterion; and source-domain change stays bounded with no class collapse.

The direct comparison between encoder adaptation and maximum head-only replay is
not statistically significant on the primary cohort, so the honest conclusion is
a \emph{replay--plasticity allocation frontier} rather than ``depth beats
volume.'' Encoder plasticity is more consistent and more replay-efficient at the
same budget, and it carries a higher compute cost; it is not proven strictly
superior to maximum replay, which is simply not established as uniquely optimal.
The frontier holds up under stress, with one caveat we state rather than bury:
the adaptation gain survives when adaptation starts from two class-reweighted
DS1 source checkpoints (of five trained, with S sensitivity $0.78$--$0.88$),
though the encoder-minus-head margins there are small and carry no reported
uncertainty; and the gain reproduces at subject level on two external databases,
where INCART ($6/6$ subjects) and SVDB ($45$--$46/46$) show $+0.14$ to $+0.24$
mean gains, with encoder LoRA again showing the higher mean than head-only
replay. Holding back a
1\,KiB firmware reserve preserves every gain over the frozen model but reshuffles
the arm ordering, which is what makes rank-1 the more robust operating point. What remains is well defined: adaptation from the three
remaining trained source variants, three RR-architecture / CNN variants that were
never trained, the unrun optimizer and higher-rank cells of the budget grid, a
larger external cohort than six INCART subjects, and---for the strong
``plasticity beats volume'' claim---a larger patient cohort that could resolve the
direct comparison, which at $n = 21$ is bounded but inconclusive (neither
superior after correction nor equivalent within the pre-registered
$\delta = 0.02$ margin; minimum detectable effect $\approx 0.01$--$0.04$).
We make no claim of microcontroller deployment; the budget is
analytical persistent state, and peak transient state is reported analytically
under three regimes (Section~\ref{sec:method}) while peak SRAM, latency, and
energy remain for future hardware-in-the-loop work.
